\documentclass[9pt]{article}

\usepackage[margin=0.5in,landscape]{geometry}
\usepackage{multicol}
\usepackage{amsmath,amssymb,amsfonts}
\usepackage{array}
\usepackage{enumitem}
\setlist{noitemsep}

\begin{document}

\begin{center}
{\Large \textbf{CSCE 222 Exam 1 Comprehensive Cheatsheet}}
\end{center}

\begin{multicols}{2}

%%%%%%%%%%%%%%%%%%%%%%%%%%%%%%%%%%%%%%%%%%%%%%%%%%%%
\section*{Logic}

\textbf{Proposition}

A declarative statement that is either true or false.

Example variables:

\[
p,q,r
\]

Truth values:

\[
T, F
\]

\subsection*{Logical Operators}

\begin{tabular}{c|c}
Symbol & Meaning \\
\hline
$\neg p$ & NOT \\
$p \land q$ & AND \\
$p \lor q$ & OR \\
$p \rightarrow q$ & Implication \\
$p \leftrightarrow q$ & Biconditional \\
$p \oplus q$ & XOR
\end{tabular}

\subsection*{Truth Tables}

\textbf{Conjunction}

\[
p \land q
\]

\begin{tabular}{c c | c}
p & q & p$\land$q\\
\hline
T&T&T\\
T&F&F\\
F&T&F\\
F&F&F
\end{tabular}

\textbf{Disjunction}

\[
p \lor q
\]

\begin{tabular}{c c | c}
p & q & p$\lor$q\\
\hline
T&T&T\\
T&F&T\\
F&T&T\\
F&F&F
\end{tabular}

\textbf{Implication}

\[
p \rightarrow q
\]

False only when:

\[
p=T, q=F
\]

Equivalent form:

\[
p\rightarrow q \equiv \neg p \lor q
\]

%%%%%%%%%%%%%%%%%%%%%%%%%%%%%%%%%%%%%%%%%%%%%%%%%%%%
\section*{Logical Equivalences}

Identity laws

\[
p \land T = p
\]

\[
p \lor F = p
\]

Domination laws

\[
p \lor T = T
\]

\[
p \land F = F
\]

Idempotent laws

\[
p \lor p = p
\]

\[
p \land p = p
\]

Double negation

\[
\neg(\neg p)=p
\]

Commutative

\[
p\land q=q\land p
\]

Associative

\[
(p\land q)\land r=p\land(q\land r)
\]

Distributive

\[
p\land(q\lor r)=(p\land q)\lor(p\land r)
\]

DeMorgan

\[
\neg(p\land q)=\neg p\lor\neg q
\]

\[
\neg(p\lor q)=\neg p\land\neg q
\]

%%%%%%%%%%%%%%%%%%%%%%%%%%%%%%%%%%%%%%%%%%%%%%%%%%%%
\section*{Rules of Inference}

\begin{tabular}{|c|c|c|}
\hline
Rule & Tautology & Name\\
\hline

$\frac{p \quad p\rightarrow q}{\therefore q}$ &
$(p\land(p\rightarrow q))\rightarrow q$
& Modus Ponens \\

\hline

$\frac{\neg q \quad p\rightarrow q}{\therefore \neg p}$ &
$(\neg q\land(p\rightarrow q))\rightarrow \neg p$
& Modus Tollens \\

\hline

$\frac{p\rightarrow q \quad q\rightarrow r}{\therefore p\rightarrow r}$ &
$((p\rightarrow q)\land(q\rightarrow r))\rightarrow(p\rightarrow r)$
& Hypothetical syllogism \\

\hline

$\frac{p\lor q \quad \neg p}{\therefore q}$ &
$((p\lor q)\land \neg p)\rightarrow q$
& Disjunctive syllogism \\

\hline

$\frac{p}{\therefore p\lor q}$ &
$p\rightarrow(p\lor q)$
& Addition \\

\hline

$\frac{p\land q}{\therefore p}$ &
$(p\land q)\rightarrow p$
& Simplification \\

\hline

$\frac{p \quad q}{\therefore p\land q}$ &
$(p\land q)\rightarrow(p\land q)$
& Conjunction \\

\hline

$\frac{p\lor q \quad \neg p\lor r}{\therefore q\lor r}$ &
$((p\lor q)\land(\neg p\lor r))\rightarrow(q\lor r)$
& Resolution \\

\hline
\end{tabular}

%%%%%%%%%%%%%%%%%%%%%%%%%%%%%%%%%%%%%%%%%%%%%%%%%%%%
\section*{Logical Errors / Fallacies}

\begin{tabular}{|c|c|c|}
\hline
Fallacy & Form & Why Wrong \\
\hline

Converse Error &
$p\rightarrow q$

$q$

$\therefore p$
&
q may occur from other causes
\\

\hline

Inverse Error &
$p\rightarrow q$

$\neg p$

$\therefore \neg q$
&
q may still be true
\\

\hline

Affirming the Consequent &
$p\rightarrow q$

$q$

$\therefore p$
&
invalid inference
\\

\hline

Denying the Antecedent &
$p\rightarrow q$

$\neg p$

$\therefore \neg q$
&
invalid inference
\\

\hline
\end{tabular}

%%%%%%%%%%%%%%%%%%%%%%%%%%%%%%%%%%%%%%%%%%%%%%%%%%%%
\section*{Predicate Logic}

Predicate:

\[
P(x)
\]

Example

\[
P(x): x>3
\]

\subsection*{Quantifiers}

Universal

\[
\forall x P(x)
\]

Existential

\[
\exists x P(x)
\]

\subsection*{Negating Quantifiers}

\[
\neg(\forall x P(x))\equiv \exists x \neg P(x)
\]

\[
\neg(\exists x P(x))\equiv \forall x \neg P(x)
\]

%%%%%%%%%%%%%%%%%%%%%%%%%%%%%%%%%%%%%%%%%%%%%%%%%%%%
\section*{Proof Techniques}

\subsection*{Direct Proof}

Goal

\[
P\rightarrow Q
\]

Steps

1 Assume $P$

2 Derive $Q$

\subsection*{Contrapositive}

\[
P\rightarrow Q
\]

prove

\[
\neg Q\rightarrow \neg P
\]

%%%%%%%%%%%%%%%%%%%%%%%%%%%%%%%%%%%%%%%%%%%%%%%%%%%%
\subsection*{Proof by Contradiction Example}

Claim:

\[
\sqrt2 \text{ is irrational}
\]

Assume opposite

\[
\sqrt2=\frac ab
\]

where

\[
a,b\in \mathbb{Z},\ gcd(a,b)=1
\]

Then

\[
2b^2=a^2
\]

So $a$ is even

Let

\[
a=2k
\]

Substitute

\[
2b^2=4k^2
\]

\[
b^2=2k^2
\]

Thus $b$ even

Contradiction since $a,b$ cannot both be even.

%%%%%%%%%%%%%%%%%%%%%%%%%%%%%%%%%%%%%%%%%%%%%%%%%%%%
\subsection*{Proof by Cases Example}

Prove

\[
n^2+n \text{ is even}
\]

Case 1

\[
n=2k
\]

\[
n^2+n=4k^2+2k=2(2k^2+k)
\]

Case 2

\[
n=2k+1
\]

\[
n^2+n=4k^2+6k+2=2(2k^2+3k+1)
\]

Both cases even.

%%%%%%%%%%%%%%%%%%%%%%%%%%%%%%%%%%%%%%%%%%%%%%%%%%%%
\section*{Even / Odd Definitions}

Even

\[
n=2k
\]

Odd

\[
n=2k+1
\]

for

\[
k\in\mathbb{Z}
\]

%%%%%%%%%%%%%%%%%%%%%%%%%%%%%%%%%%%%%%%%%%%%%%%%%%%%
\section*{Rational / Irrational}

Rational

\[
x=\frac ab
\]

where

\[
a,b\in\mathbb Z,\ b\neq0
\]

Irrational

cannot be expressed as ratio of integers.

%%%%%%%%%%%%%%%%%%%%%%%%%%%%%%%%%%%%%%%%%%%%%%%%%%%%
\section*{Sets}

Subset

\[
A\subseteq B
\]

iff

\[
\forall x(x\in A\rightarrow x\in B)
\]

Set operations

Union

\[
A\cup B
\]

Intersection

\[
A\cap B
\]

Difference

\[
A-B
\]

%%%%%%%%%%%%%%%%%%%%%%%%%%%%%%%%%%%%%%%%%%%%%%%%%%%%
\section*{Set Identities}

\begin{tabular}{|c|c|}
\hline
Identity & Name\\
\hline

$A\cap U=A$

$A\cup\emptyset=A$
&
Identity Laws
\\

\hline

$A\cup U=U$

$A\cap\emptyset=\emptyset$
&
Domination Laws
\\

\hline

$A\cup A=A$

$A\cap A=A$
&
Idempotent
\\

\hline

$\overline{\overline A}=A$
&
Complementation
\\

\hline

$A\cup B=B\cup A$

$A\cap B=B\cap A$
&
Commutative
\\

\hline

$A\cup(B\cup C)=(A\cup B)\cup C$

$A\cap(B\cap C)=(A\cap B)\cap C$
&
Associative
\\

\hline

$A\cup(B\cap C)=(A\cup B)\cap(A\cup C)$

$A\cap(B\cup C)=(A\cap B)\cup(A\cap C)$
&
Distributive
\\

\hline

$\overline{A\cap B}=\overline A\cup\overline B$

$\overline{A\cup B}=\overline A\cap\overline B$
&
DeMorgan
\\

\hline

$A\cup(A\cap B)=A$

$A\cap(A\cup B)=A$
&
Absorption
\\

\hline

$A\cup\overline A=U$

$A\cap\overline A=\emptyset$
&
Complement
\\

\hline
\end{tabular}

%%%%%%%%%%%%%%%%%%%%%%%%%%%%%%%%%%%%%%%%%%%%%%%%%%%%
\section*{Cardinality}

Cardinality

\[
|A|
\]

= number of elements.

Power set

\[
P(A)
\]

Number of subsets

\[
|P(A)|=2^{|A|}
\]

Cartesian product

\[
A\times B
\]

\[
|A\times B|=|A||B|
\]

%%%%%%%%%%%%%%%%%%%%%%%%%%%%%%%%%%%%%%%%%%%%%%%%%%%%
\section*{Functions}

Function

\[
f:A\rightarrow B
\]

Injective

\[
f(x_1)=f(x_2)\Rightarrow x_1=x_2
\]

Surjective

\[
\forall y\in B,\exists x\in A:f(x)=y
\]

Bijective

injective + surjective

%%%%%%%%%%%%%%%%%%%%%%%%%%%%%%%%%%%%%%%%%%%%%%%%%%%%
\section*{Algorithms and Complexity}

Common complexities

\begin{tabular}{c|c}
Type & Growth\\
\hline
Constant & $O(1)$\\
Logarithmic & $O(\log n)$\\
Linear & $O(n)$\\
Quadratic & $O(n^2)$\\
Polynomial & $O(n^k)$\\
Exponential & $O(2^n)$
\end{tabular}

%%%%%%%%%%%%%%%%%%%%%%%%%%%%%%%%%%%%%%%%%%%%%%%%%%%%
\section*{Asymptotic Notation}

Big-O

\[
f(n)\le Cg(n)
\]

Big-Omega

\[
f(n)\ge Cg(n)
\]

Big-Theta

\[
Cg(n)\le f(n)\le C_2g(n)
\]

%%%%%%%%%%%%%%%%%%%%%%%%%%%%%%%%%%%%%%%%%%%%%%%%%%%%

\end{multicols}

\end{document}